% ৩.২ সংখ্যা পদ্ধতি
\section{সংখ্যা পদ্ধতি}

\subsection{সংখ্যা পদ্ধতির ধারণা}
\begin{itemize}
\item সংখ্যা পদ্ধতি কী, অঙ্ক/ Digit
\item ভিত্তি / Base বা Radix
\item Place value ও Face value
\item Positional ও Non-positional system
\end{itemize}

\subsection{সংখ্যা পদ্ধতির প্রকারভেদ}
\paragraph{Decimal (Base 10)}\mbox{}\\
\begin{itemize}\item Digit: 0--9\item দৈনন্দিন ব্যবহার\end{itemize}

\paragraph{Binary (Base 2)}\mbox{}\\
\begin{itemize}\item Digit: 0,1\item কম্পিউটারের মূল সংখ্যা পদ্ধতি\end{itemize}

\paragraph{Octal (Base 8) and Hexadecimal (Base 16)}\mbox{}\\
\begin{itemize}\item Octal: 0--7\item Hex: 0--9, A--F\end{itemize}
